% !PRINTED none
\begin{center}
\gesperrt{Twenty-fifth Chapter.}

Functionals.

\rule{4cm}{0.4pt}
\end{center}

\begin{center}
§ 177. \gesperrt{Rational Functionals.}
\end{center}

We now carry over the concept of the functional, which served us in
the 17th chapter of the second volume for the foundation of the theory
of algebraic numbers, to the algebraic functions.

We therefore adjoin arbitrary variables to the field $\Omega$ and
calculate with them according to the rules of literal calculation. There
thus arises an extended field $\overline{\Omega}$, whose elements are called \gesperrt{functionals}.
 The field $\overline{\Omega}$ itself is called the \gesperrt{functional field}.
 The variables here are not to be taken, in the sense of analysis,
as symbols for variable numbers, but as mere symbols of calculation,
and are to be well distinguished from the
variables $z, \theta$ of the field $\Omega$. We wish to call these auxiliary variables
the \gesperrt{functional variables}.

From the field $Z$ of the rational functions of $z$ there arises,
by adjunction of the variables, the \gesperrt{field $\overline{Z}$ of the rational
functionals}.

An integral rational function of the variables:

\[
\Phi(u, v, w, \ldots)
\]

with integral rational functions of $z$ as coefficients is called
an integral function in $Z$. The greatest common divisor
of the coefficients of $\Phi$ is called \gesperrt{the divisor of the function}, and
the function is called \gesperrt{primitive} if the coefficients have no common
 divisor. The primitive rational functions
 and their quotients are also called \gesperrt{units in the field}
$\overline{Z}$.
